\section{Geometric transfer, the positive adjoint, and the exact
reflection-average
defect}\label{geometric-transfer-the-positive-adjoint-and-the-exact-reflection-average-defect}

24 September 2026. Proof locators GTAH0--GTAH7. This carries out the
next attempt requested by the user's rule: after identifying what a
result has not established, ask what would establish it using the entire
known programme, and try that construction.

\subsection{GTAH0. The actual objects already
constructed}\label{gtah0.-the-actual-objects-already-constructed}

All arithmetic and coefficient operations below follow the complete
arithmetic reconstruction. The supporting datum remains
\(\tau\langle Z_1;\text{no intrinsic }Z_2\rangle\); none of these
operations defines addition, coordinates, or a distance on that datum.
The complete source quotient \(\mathcal Q=\mathcal B/\mathcal I\), its
full multiplicity jets and all its original-zeta factors remain as in
SSI, RD, GCF and PGF. This calculation does not replace that quotient by
its value observation.

GIQ2 proves the continuous value map with dense image \[
E_0:\mathcal Q\longrightarrow\mathcal H,
\qquad E_0[F]=(F(\rho))_{\rho\in\mathscr Z},
\qquad
\mathcal H=\ell^2(\mathscr Z,m_\rho),
\tag{GTAH0.1}
\] where \(\mathscr Z\) is the entire set of distinct nontrivial zeros
of the original \(\zeta\), and \(m_\rho\) is its actual multiplicity.
Its kernel is exactly the ideal of classes with zero values at every
zero; higher jets remain in this kernel and in the source quotient. The
topology here is the stated positive value norm, not a new topology
silently assigned to \(\mathcal Q\). RD1 supplies explicit
representatives with a prescribed full jet at one zero and every
required vanishing jet at all other zeros. Thus finite value vectors
actually lie in the image, which also proves its density after the SSI
identification.

Use the inner product linear in its first variable, \[
\langle x,y\rangle=\sum_\rho m_\rho x_\rho\overline{y_\rho}.
\tag{GTAH0.2}
\] The original functional equation gives the involution
\(\rho^\#=1-\overline\rho\), with \(m_{\rho^\#}=m_\rho\). The operator
\[
(\mathsf Jx)_\rho=x_{\rho^\#},\qquad
\mathsf J^2=I,\quad\mathsf J^*=\mathsf J
\tag{GTAH0.3}
\] is a unitary involution on this complete value space. The full
original Weil form, including its original arithmetic expression in
GIQ9, is \[
W(x,y)=\langle x,\mathsf Jy\rangle
=\sum_\rho m_\rho x_\rho\overline{y_{\rho^\#}}.
\tag{GTAH0.4}
\] GIQ9 proves the receiving map from the original tests and the
complete prime, Gamma, endpoint and finite trivial-zero formulas. Only
the inner-product convention changes here: GIQ uses conjugate linearity
in the first variable; (GTAH0.4) is its complex conjugate. No sign
changes in quadratic values result.

The arithmetic dilation and the actual degree-weighted transfer on this
receiver are \[
(\mathsf T_a x)_\rho=a^\rho x_\rho,
\qquad
\mathsf U_a=a\mathsf T_{1/a},\qquad a>0.
\tag{GTAH0.5}
\] For recovered positive integers \(a=n\), the factor \(n\) is the
actual cover degree proved in GTR0--GTR11, including the multiplicities
at both poles. For general \(a>0\) these are coefficient operators; no
nonintegral sphere cover is asserted. The known strip \(0<\Re\rho<1\)
makes both operators bounded, with bounded inverses. Direct
multiplication gives \[
\mathsf U_a\mathsf T_a=\mathsf T_a\mathsf U_a=aI,
\qquad \mathsf T_a\mathsf T_b=\mathsf T_{ab},
\qquad \mathsf U_a\mathsf U_b=\mathsf U_{ab}.
\tag{GTAH0.6}
\] The original coefficient maps commute with \(E_0\). These statements
hold on the full actual zero set, without a zero-height cutoff or
simplicity assumption.

\subsection{GTAH1. What would give the desired modulus, and the direct
calculation}\label{gtah1.-what-would-give-the-desired-modulus-and-the-direct-calculation}

The existing geometric degree identity would give the common modulus if
its transfer were the adjoint for the actual positive value norm. The
direct calculation of that adjoint is \[
(\mathsf T_a^*x)_\rho=a^{\overline\rho}x_\rho,
\qquad
\mathsf J\mathsf U_a\mathsf J=\mathsf T_a^*.
\tag{GTAH1.1}
\] Indeed substitution in (GTAH0.2) proves the first equality. For the
second, the multiplier on the left is
\(a^{1-\rho^\#}=a^{\overline\rho}\). No use of a conjectural real part
is made.

For the actual Weil form the transfer already is the adjoint: \[
W(\mathsf T_a x,y)=W(x,\mathsf U_a y),
\qquad W(\mathsf T_a x,\mathsf T_a y)=aW(x,y).
\tag{GTAH1.2}
\] Both follow directly from \(\overline{\rho^\#}=1-\rho\). Thus the
precise discrepancy between the two adjoints is the bounded operator \[
\mathsf D_a:=\mathsf T_a^*-\mathsf U_a,
\qquad
(\mathsf D_a x)_\rho
=\bigl(a^{\overline\rho}-a^{1-\rho}\bigr)x_\rho.
\tag{GTAH1.3}
\] For \(\rho=\sigma+i\gamma\), all scalar factors remain explicit: \[
a^{\overline\rho}-a^{1-\rho}
=e^{-i\gamma\log a}\bigl(a^\sigma-a^{1-\sigma}\bigr).
\tag{GTAH1.4}
\] For any fixed \(a>1\), its kernel is exactly the closed subspace
supported on the actual critical-line zeros. This is proved by strict
monotonicity of the real function \(a^x\), applied coordinatewise. In
particular, every possible off-line value is retained by this defect,
whereas every higher jet lies in the separately retained kernel of
\(E_0\). This is a computation of the candidate positive adjoint, not a
conclusion about which off-line zeros exist.

\subsection{GTAH2. Reflection averaging is a concrete second
attempt}\label{gtah2.-reflection-averaging-is-a-concrete-second-attempt}

The positive form is obtained from the existing Weil form by
\(\langle x,y\rangle=W(x,\mathsf Jy)\). Rather than declare this change
harmless, retain \(\mathsf J\) and test the reflected average of the
actual operators: \[
\widetilde{\mathsf T}_a
=\frac12\bigl(\mathsf T_a+\mathsf J\mathsf T_a\mathsf J\bigr),
\qquad
\widetilde{\mathsf U}_a
=\frac12\bigl(\mathsf U_a+\mathsf J\mathsf U_a\mathsf J\bigr).
\tag{GTAH2.1}
\] These are bounded operators on the entire value space, and \[
\widetilde{\mathsf T}_a^*=\widetilde{\mathsf U}_a.
\tag{GTAH2.2}
\] Proof: transpose (GTAH2.1), use (GTAH1.1), and use
\(\mathsf J^*=\mathsf J\). Thus this attempt does repair the
positive-adjoint identity. It is not asserted that \(\mathsf J\) or
these averages lift to the original test quotient; their domain is the
proved value completion.

The complete scalar multiplier of the first average is \[
\widetilde t_a(\rho)
=\frac{a^\rho+a^{\rho^\#}}2
=e^{i\gamma\log a}\frac{a^\sigma+a^{1-\sigma}}2.
\tag{GTAH2.3}
\] Its degree relation is consequently \[
\boxed{\widetilde{\mathsf U}_a\widetilde{\mathsf T}_a
=aI+\frac14\mathsf D_a^*\mathsf D_a.}
\tag{GTAH2.4}
\] To prove this identity, square the absolute value in (GTAH2.3),
retaining the cross term \(2a\), and subtract \(a\): \[
\left(\frac{a^\sigma+a^{1-\sigma}}2\right)^2-a
=\frac{(a^\sigma-a^{1-\sigma})^2}{4}.
\tag{GTAH2.5}
\] This is exactly the diagonal multiplier of the extra operator in
(GTAH2.4). Hence the attempted positive repair produces an explicitly
positive degree defect rather than making the original degree identity
automatic.

\subsection{GTAH3. All arithmetic compositions and the return-time
defect}\label{gtah3.-all-arithmetic-compositions-and-the-return-time-defect}

The same defect also measures failure of the averaged operators to
preserve the full multiplicative arithmetic action. Put \[
\mathsf E_a=\mathsf T_a-\mathsf J\mathsf T_a\mathsf J.
\tag{GTAH3.1}
\] All these operators are diagonal, so they commute. Expanding the four
terms gives, for every \(a,b>0\), \[
\boxed{\widetilde{\mathsf T}_a\widetilde{\mathsf T}_b
-\widetilde{\mathsf T}_{ab}
=-\frac14\mathsf E_a\mathsf E_b.}
\tag{GTAH3.2}
\] The proof uses both original group laws in (GTAH0.6) and the same
laws after conjugation by \(\mathsf J\); it makes no selection of two
named primes. The right side is a global formula for every recovered
arithmetic parameter and all its repetitions.

Taking \(b=1/a\) in the exact scalar multipliers yields \[
\boxed{\widetilde{\mathsf T}_a\widetilde{\mathsf T}_{1/a}
=I+\frac1{4a}\mathsf D_a^*\mathsf D_a.}
\tag{GTAH3.3}
\] Thus averaging is compatible with inverse return only on the kernel
already calculated in GTAH1. The calculation connects the positive
repair to the complete return action, rather than treating a changed
quadratic form as if all earlier maps remained fixed.

The scalar formula also gives
\(\widetilde{\mathsf U}_a=a\widetilde{\mathsf T}_{1/a}\), so (GTAH3.3)
recovers (GTAH2.4) with the same degree. Both defects are the same
operator with the explicitly retained factor \(a\).

\subsection{GTAH4. A single bounded self-adjoint operator controls the
whole
defect}\label{gtah4.-a-single-bounded-self-adjoint-operator-controls-the-whole-defect}

On the value space let the full generator \(\mathsf L\) have its maximal
diagonal domain \[
\operatorname{Dom}\mathsf L
=\{x:\sum_\rho m_\rho|\rho|^2|x_\rho|^2<\infty\},
\qquad (\mathsf Lx)_\rho=\rho x_\rho.
\tag{GTAH4.1}
\] Finite value vectors form a dense domain. Direct use of (GTAH0.2)
shows that its adjoint has the same domain and multiplier
\(\overline\rho\). The operator \(\mathsf L+\mathsf L^*-I\) on that
domain has the bounded extension with multiplier \(2\Re\rho-1\). Define
its half by the actual comparison \[
\mathsf C=\frac12(\mathsf L+\mathsf L^*-I),\qquad
(\mathsf Cx)_\rho=(\sigma-\tfrac12)x_\rho,
\quad \|\mathsf C\|\le\tfrac12,
\quad \mathsf J\mathsf C\mathsf J=-\mathsf C.
\tag{GTAH4.2}
\] This notation records all terms of the generator comparison; the
scalar \(1/2\) is not a coordinate assigned to \(\tau\). It does not
replace \(\mathsf L\), \(\zeta\), or their original factors.

The exact identity for every \(a>0\) is \[
\mathsf D_a^*\mathsf D_a
=4a\sinh^2((\log a)\mathsf C).
\tag{GTAH4.3}
\] It follows by writing \(a^\sigma-a^{1-\sigma}
=2a^{1/2}\sinh((\sigma-1/2)\log a)\) in (GTAH1.4), keeping both powers
in the preceding expression. Bounded functional calculus here means
simply applying the indicated continuous function to the diagonal real
coordinates, so no spectral-domain assertion is concealed.

For \(a>1\), the complete operator estimates are \[
4a(\log a)^2\mathsf C^2
\ \le\ \mathsf D_a^*\mathsf D_a
\ \le\ 4(a-1)^2\mathsf C^2.
\tag{GTAH4.4}
\] Indeed for \(0\le x\le1/2\), \(\sinh((\log a)x)\ge(\log a)x\) by its
power series, while convexity and its zero value at zero give
\(\sinh((\log a)x)\le 2x\sinh((\log a)/2)\). Squaring, multiplying by
\(4a\), and using \(4a\sinh^2((\log a)/2)=(a-1)^2\) proves the bounds.
Evenness handles negative \(x\). These are global analytic estimates,
not numerical bounds on a finite list of zeros.

\subsection{GTAH5. The exact positive trace with all
multiplicities}\label{gtah5.-the-exact-positive-trace-with-all-multiplicities}

The ordinary operator trace on \(\ell^2(\mathscr Z,m_\rho)\) does not
itself insert multiplicity \(m_\rho\): rescaling an orthonormal basis
would remove the weight. To preserve the original zeta multiplicities,
use the actual diagonal trace \[
\operatorname{tr}_\zeta(M_f)=\sum_\rho m_\rho f(\rho)
\tag{GTAH5.1}
\] on the diagonal algebra, first for nonnegative \(f\), and then when
the absolute weighted sum converges. This is the trace received from the
full local multiplication algebra: multiplication by a germ
\(f(\rho+h)\) on \(\mathbb C[h]/h^{m_\rho}\) has triangular matrix with
all \(m_\rho\) diagonal entries \(f(\rho)\). Its trace is therefore
\(m_\rho f(\rho)\), also equal to
\(\operatorname{Res}_\rho f\,\zeta'/\zeta\). The nilpotent terms remain
in that matrix, although their trace is zero. This is the established
PCJ10/GIQ trace, with its original-zeta receiver.

For every \(t>0\), set \(k_t(\rho)=e^{-t(\Im\rho)^2}\). The
unconditional full zero-count estimate
\(\sum_{|\rho|\le R}m_\rho=O(R\log(R+2))\) proves
\(\sum m_\rho k_t(\rho)<\infty\): bound each strip \(j\le|\Im\rho|<j+1\)
by its containing disk count and sum the convergent series
\(O((j+2)\log(j+3)e^{-tj^2})\). Since \(|a^\sigma-a^{1-\sigma}|\le a-1\)
for \(a>1\), the entire positive defect trace is finite and exactly \[
\mathcal E(a,t):=
\operatorname{tr}_\zeta(M_{k_t}\mathsf D_a^*\mathsf D_a)
=\sum_\rho m_\rho e^{-t(\Im\rho)^2}
\bigl(a^{\Re\rho}-a^{1-\Re\rho}\bigr)^2.
\tag{GTAH5.2}
\] No zero term is cut off. Reflection preserves every summand, since it
negates the difference before squaring and leaves \(\Im\rho\) unchanged.
Thus the two reflected directions add to this positive quantity; their
signed first differences cancel, but their exact degree defect does not.

Taking this faithful positive diagonal trace in (GTAH4.4) gives \[
4a(\log a)^2\sum_\rho m_\rho k_t(\rho)(\Re\rho-\tfrac12)^2
\le\mathcal E(a,t)
\le4(a-1)^2\sum_\rho m_\rho k_t(\rho)(\Re\rho-\tfrac12)^2.
\tag{GTAH5.3}
\] The trace is a receiver on actual spectral values; no interchange of
a nonholomorphic test with the holomorphic prime explicit formula has
been asserted. Its vanishing has not been presumed. The purpose of
(GTAH5.2) is to retain and measure the exact error in the attempted
construction, not to rename a vanishing conjecture as a theorem.

\subsection{GTAH6. The result of the attempted repair and its next
mathematical
receiver}\label{gtah6.-the-result-of-the-attempted-repair-and-its-next-mathematical-receiver}

The attempt began with the actual geometric transfer and the positive
norm already constructed in GIQ. Its adjoint was calculated, not
assigned. Reflection averaging repairs that adjoint identity, but the
original degree and arithmetic composition laws acquire the same
positive defect (GTAH2.4) and (GTAH3.3). The full scalar, operator and
multiplicity-trace descriptions of that defect are (GTAH1.4), (GTAH4.3)
and (GTAH5.2).

The next useful target from this finding and the complete programme is
the original source-to-dual map, before its value quotient: can the
original-zeta residue duality provide an adjoint comparison whose image
forces this defect to vanish? That attempt is now carried out in
\nolinkurl{RESIDUE_TRACE_TRANSFER_INDEPENDENT.md}, RTT0--RTT11, read in
full for this receiving calculation. It constructs the actual continuous
anti-linear map \[
\mathsf P=L^{-2}M_{s^2\zeta'(1-s)}\mathsf K:\mathcal Q\to\mathcal Q,
\qquad (\mathsf KG)(s)=\overline{G(\overline s)},
\tag{GTAH6.1}
\] and the full strong residue map
\(\mathsf S:\mathcal Q\to\mathcal H_{\rm res}\) with \[
\widehat{\mathcal R}(F,\mathsf SG)=W(E_0F,E_0G),
\qquad \ker\mathsf S=\ker E_0=\mathcal N_0.
\tag{GTAH6.2}
\] The notation \(\mathsf K\) here is RTT's conjugation, distinct from
the self-adjoint operator \(\mathsf C\) in GTAH4. The multiplier has the
exact endpoint value \(s^2\zeta'(1-s)|_{s=0}=-1\); the inverse of \(L\)
is proved on \(\mathcal Q\), rather than division through an unremoved
source pole. RTT3 calculates rank one on every full multiplicity block,
with its complete nonzero \(\zeta'\) coefficient. RTT5 proves the strong
limit on the entire quotient without a finite-support density
assumption.

This map permits a further exact test on the original source. Define the
positive continuous form \[
\mathfrak d_a(F,G)=\langle\mathsf D_aE_0F,\mathsf D_aE_0G\rangle.
\tag{GTAH6.3}
\] Continuity follows from continuity of \(E_0\) and boundedness of
\(\mathsf D_a\). Its kernel for \(a>1\) consists exactly of classes
whose value is zero at every off-line zero; the higher-jet kernel
\(\mathcal N_0\) is contained in it. Since (GTAH6.2) identifies exactly
that higher-jet kernel, (GTAH6.3) descends without further loss to the
trace image: give \(\mathsf S(\mathcal Q)\) the quotient topology from
\(\mathcal Q/\mathcal N_0\), and define \[
\mathfrak d_{a,\rm res}(\mathsf SF,\mathsf SG)
=\overline{\mathfrak d_a(F,G)}.
\tag{GTAH6.4}
\] The conjugation is required because \(\mathsf S\) is anti-linear; it
makes the displayed form linear in its first argument on the target.
Changing representatives by \(\mathcal N_0\) leaves it unchanged. This
proves a positive continuous target form for the stated quotient
topology, not for an unverified subspace topology of
\(\mathcal H_{\rm res}\). Its quadratic value is exactly (GTAH6.3), so
no off-line value is erased by this source-to-residue operation.

The signed test is equally explicit. On any two-element orbit
\(\{\rho,\rho^\#\}\) of the actual zero set, use the actual global
isolators and put \(F=e_{\rho,0}-e_{\rho^\#,0}\). Then \[
W(E_0F,E_0F)=-2m_\rho,
\qquad
\widehat{\mathcal R}(F,\mathsf SF)=-2m_\rho,
\tag{GTAH6.5}
\] by its two nonzero values and (GTAH6.2). This quantifies over the
actual orbits; it does not assert that any off-line orbit exists. It
proves that the infinitesimal kernel in this particular map cannot
secretly remove such a direction while retaining the original Weil
pairing.

There is also a direct comparison for the entire value Hilbert space,
rather than only the quotient image just used. RTT6.2 constructs \[
\mathsf A:H\longrightarrow\mathcal H_{\rm res},\qquad
(\widehat\iota\mathsf A y)(F)=\langle E_0F,Jy\rangle,
\qquad \mathsf S=\mathsf A E_0.
\tag{GTAH6.6}
\] For a bounded subset \(B\subset\mathcal Q\), its defining strong
seminorm is at most \(\sup_{F\in B}\|E_0F\|\,\|y\|\), so \(\mathsf A\)
is continuous and anti-linear. If \(\mathsf Ay=0\), then \(Jy\) is
orthogonal to the dense subspace \(E_0(\mathcal Q)\); hence \(Jy=0\) and
\(y=0\). Applying this injectivity to each vector gives the exact
equivalence \[
\mathsf A\mathsf D_a^*\mathsf D_a=0
\quad\Longleftrightarrow\quad
\mathsf D_a^*\mathsf D_a=0
\quad\Longleftrightarrow\quad
\mathsf D_a=0.
\tag{GTAH6.7}
\] The last equivalence follows from
\(\langle\mathsf D_a^*\mathsf D_ax,x\rangle=\|\mathsf D_ax\|^2\) for
every \(x\in H\). Thus the complete residue receiver detects the
positive degree defect faithfully. This does not assert that
\(\mathsf A(H)\) has a continuous inverse with its subspace topology, or
that the nonholomorphic value multiplier lifts to \(\mathcal Q\).

Thus both attempts are calculated: reflection averaging and the full
source-to-residue trace map. The first yields the exact degree defect;
the second carries that defect, with all value information, into the
specified residue receiver. No conclusion about absence of off-line
zeros or the full target's Deligne weights has been assumed. The next
source-level construction must now control this explicit surviving form
while preserving the original transfer and complete return action; the
linked full strip-multiplier calculation is addressing the source
modules rather than substituting another scalar criterion.

\subsection{GTAH7. Exact source and proof
provenance}\label{gtah7.-exact-source-and-proof-provenance}

Programme proofs used: GIQ2,3,8,9 in
\nolinkurl{GLOBAL_INFINITESIMAL_QUOTIENT.md}; RD1--RD5 in
\nolinkurl{ORIGINAL_ZETA_RESIDUE_DUAL_AND_SUPPORT_MAPS.md}; GTR0--GTR11
in \nolinkurl{CC_RAMIFIED_TRACE_AND_RESIDUE_ACTION.md}; SSI0--SSI8 in
\nolinkurl{EXACT_SCHWARTZ_SUMMATION_IMAGE.md}; PCJ10 in
\nolinkurl{POSITIVE_COMPLETION_AND_ZERO_JETS.md}. These inputs already
prove the underlying zero/value completion, trace radical, actual degree
transfer, full source ideal and all-zero pairing. The new calculation is
the exact adjoint-repair attempt with the full transfer composition and
its positive defect; it does not claim the earlier constructions as new.

Human-source route: Connes and Consani, \emph{Schemes over F1 and zeta
functions}, \href{https://arxiv.org/abs/0903.2024}{arXiv:0903.2024v3},
§5, for the original summation/cohomology construction; Deligne,
\emph{La conjecture de Weil. II},
\href{https://www.numdam.org/item/PMIHES_1980__52__137_0/}{IHÉS52(1980),137--252},
§§1.5 and3.6, for the distinct use of positivity and weight separation.
No finite-rank étale realization or Deligne weight theorem is assumed
for the spaces in this note.
